\section{Replacing indicators by arithmetic models}\label{sec:comparison}

The comparison argument depends on how an indicator interacts with
selected primes and on one-variable distribution, rather than on a
particular Euler product. We isolate these requirements. The arithmetic
applications below construct the models and check the requirements;
the signed graph comparison and its frequency estimates are proved once.

Fix an even integer \(L_0\ge2\) and a reduced class
\(c\in(\Z/L_0\Z)^\times\). At the outer scale \(X\), write
\[
 L=\log X,\qquad y=\exp(L^{2/5}),\qquad
 S=\{p:p\le y\},\qquad E=\exp(\sqrt L).
\]
For a set of primes \(\mathcal P\), let \(n_{\mathcal P}\) be the
part of \(n\) supported on \(\mathcal P\), including its full
valuations. A rough integer has no prime factor in \(S\).
There are two families of lines in the theorem:
\begin{align*}
 \mathcal L_X^{\rm sum}&=\{(u,v)\in\Z_{>0}^2:Au+Bv=X\},
       &t(u,v)&=Au, &J_X&=[0,X],\\
 \mathcal L_X^{\rm shift}&=\{(u,u+h):u\in\Z,\ X\le u\le2X\},
       &t(u,v)&=u, &J_X&=[X,2X].
\end{align*}
In the first family \(X\) is an integer target and \(A,B\) are
fixed coprime positive integers supported on primes dividing \(L_0\).
In the second, \(h\ne0\) is a fixed integer and \(X\) tends to
infinity; both coordinates are then positive. Set \(k=X\) in the
sum family and \(k=h\) in the shift family. Dependence on the
fixed data, including \(h\) in the second family, is permitted.

Let \(0\le f_i\le1\), \(i=1,2\), be fixed functions on the
positive integers. For every sufficiently large fixed integer \(K\),
let \(D_{i,X,K}\) be an arithmetic model, and write
\(\xi_{i,X,K}=f_i-D_{i,X,K}\). We often suppress \(X,K\).
The entire model, its cutoff \(y\), and its small-prime records stay
tied to this outer scale during all subsequent dilations.

Choose a positive comparison scale \(\mathcal S_X\) satisfying
\begin{equation}\label{eq:comparison-scale}
        XL^{-\kappa_0}\le\mathcal S_X\le XL^{\kappa_0}
\end{equation}
for a fixed \(\kappa_0\ge0\). Arithmetic factors may be included
in this scale. The hypotheses are as follows.

\begin{description}
\item[\textup{(C1)} Bounded models.]
For every fixed \(K\), sufficiently large \(X\) gives
\(0\le D_{i,X,K}(n)\le1\) on \(1\le n\le X^3\).
The bound one is independent of the model depth; its threshold
may depend on \(K\).

\item[\textup{(C2)} A common prime class.]
For every prime \(p\le y\) with \(p\equiv c\pmod{L_0}\),
\[
 f_i(pn)=f_i(n),\qquad D_{i,X,K}(pn)=D_{i,X,K}(n)
\]
whenever the arguments are in the range under consideration.
This includes repeated multiplication by \(p\). A residue filter
is part of \(f_i\) and must satisfy the same equality.

\item[\textup{(C3)} Centered one-variable rough means.]
For every fixed \(B_0,D_0>0\), there is a depth \(K_0\) such that
every fixed \(K\ge K_0\) has the following property. Uniformly for
\(S\)-supported integers \(r\le E\),
\(X^{1/2}\le Y\le X^2\), bounded periodic functions \(a(v)\)
of period at most \(L^{B_0}\), and smooth functions \(W\) supported
in a fixed compact subinterval of \((0,\infty)\),
\begin{equation}\label{eq:comparison-rough-means}
 \left|\sum_{v\ {\mathrm{rough}}}
       \xi_{i,X,K}(rv)a(v)W(rv/Y)\right|
                       \ll_{B_0,D_0,K}\frac Yr L^{-D_0}.
\end{equation}
Here \(|a|\le1\), and every fixed logarithmic derivative of \(W\)
used in the estimate is at most a fixed power of \(L^{B_0}\).
The hypothesis permits Mellin factors with polynomial-logarithmic
frequency, by including them in \(W\). Its constants may depend on
the chosen derivative orders. Every surviving character branch,
including a negative or complex pole parameter, must be retained
when verifying this estimate.

\item[\textup{(C4)} A positive sparse record envelope.]
Take disjoint prime banks in the two ranges
\[
 [e^{L^{1/8}},e^{L^{1/4}}],\qquad
 [e^{L^{1/3}},e^{L^{2/5}}],
\]
each bank consisting of the primes in an interval which are
\(c\pmod{L_0}\), after deleting primes dividing \(k\). Assume each
bank has reciprocal mass between \(1/2\) and \(2\), and that the
number in either range is at most \((\log L)/(16\varphi(L_0))\).
Write \(V_g=\sum_{p\in\mathcal P_g}p^{-1}\) for its reciprocal mass.
Let \(\mathcal B\) be their union. Define the independent joint
record law \(\mu_{\mathcal B}=\prod_{p\in\mathcal B}\mu_p\) by
\begin{equation}\label{eq:bank-law}
 \mu_p(0,0)=1-2/p,\qquad
 \mu_p(e,0)=\mu_p(0,e)=(1-1/p)p^{-e}\quad(e\ge1),
\end{equation}
and zero mass on two positive valuations.
For every fixed \(B_0>0\), nonnegative record function
\(F(r,s)\le L^{B_0}\), interval \(J\subset J_X\), and
\(w_i\in\{f_i,D_{i,X,K}\}\), require
\begin{align}
 &\sum_{\substack{(u,v)\in\mathcal L_X\\t(u,v)\in J}}
     F(u_{\mathcal B},v_{\mathcal B})w_1(u)w_2(v)
       \ll_{B_0,K}\frac{|J|}{X}\mathcal S_X
                              \E_{\mu_{\mathcal B}}F\notag\\
 &\hspace{34mm}
       +X^{1/2}L^{C_{B_0}}
       +XL^{C_{B_0}}\exp(-c_0L^{1/10}).
                                                        \label{eq:comparison-envelope}
\end{align}
The positive constant \(c_0\) and exponent \(C_{B_0}\) are
independent of the final fixed model depth. Constants and thresholds
may depend on that depth. Uniformity includes the bank intervals,
the exact valuations in the record, and the interval endpoints.
\end{description}

For a bank \(g\), let \(\omega_g(n)\) count its distinct prime
divisors, and define
\[
 I_g(u,v)=\1_{\omega_g(u)=\omega_g(v)=1},\qquad
 \Omega_T(n)=\prod_{g\in T}V_g^{-1}\1_{\omega_g(n)=1}.
\]
Here \(T\) is a set of bank indices. We also write
\(\Omega_{\mathcal P}\) when \(\mathcal P\) denotes the union of
the active banks. The joint valuation law above
is the ordinary CRT law on either line: every bank prime avoids
\(k\) and is a unit at the fixed coefficients.

Condition (C3) is a statement about one variable on a progression;
it allows arbitrary retained small-prime records, but no arbitrary
weight in the rough variable. Condition (C4) is only an upper bound
with nonnegative weights. In the applications it follows from CRT
root counts and an upper sieve, with a power-smaller endpoint error.
Neither condition assumes centered binary cancellation or a major-arc
energy estimate. In particular, choosing a model with the right
average density alone does not verify (C3).

\begin{theorem}[Arithmetic model comparison]\label{thm:model-comparison}
Assume the zero-free and weighted-dilation inputs of
Section~\ref{sec:inputs}, and conditions \textup{(C1)--(C4)}.
Put \(\delta=1/(8192\varphi(L_0))\). For every sufficiently deep
fixed \(K\), uniformly for intervals \(J\subset J_X\),
\begin{equation}\label{eq:model-comparison-conclusion}
 \sum_{\substack{(u,v)\in\mathcal L_X\\t(u,v)\in J}}
     \bigl(f_1(u)f_2(v)-D_{1,X,K}(u)D_{2,X,K}(v)\bigr)
                       \ll_K\mathcal S_XL^{-\delta}.
\end{equation}
For the sum family the estimate is uniform in the integer target
and its prime valuations. For the shift family the integer \(h\)
is fixed, and the constant may depend on it. The two families have
different major-arc geometries; no uniformity in the intervening
range of growing sublinear shifts is asserted.
\end{theorem}

The proof occupies Sections~\ref{sec:graph} and
\ref{sec:fixed-shift-comparison}, followed by the quantitative
reconstruction in Section~\ref{sec:assembly}. A fixed reduced prime
class supplies the same number of banks in both cases. The graph
comparison produces two independent large prime products. For a
reflected target, its major arcs form target-sized rational grids;
recordwise centering and a small-prime extraction control their
residue energy. For a fixed shift the arcs are wider, and the existing
additive-to-Dirichlet conversion instead gives low and high Mellin
frequencies. Condition (C3) handles the low range; exact invariance,
the same zero-free prime estimate, and the ordinary Dirichlet
mean-value bound handle the high range. Both estimates are proved
from the stated hypotheses below.

Two scope distinctions are useful. A pair of subgroups always has
the identity prime class available. A self-coset prescription has
any one of its prime classes available, with a slot-dependent unit
phase in the character frame. Arbitrary pairs of prime prescriptions
can have no common invariant prime class, so they do not satisfy
(C2). Also, the comparison saving does not bound the secondary
terms of a particular Euler product: the final asymptotic rate must
also pay for evaluating its complete arithmetic model.
